\chapterNote{Based on a talk given on December 16, 2022}

\chapter{Me and Mine}
\tocChapterNote{Based on a talk given on December 16, 2022}
%\tocChapterNote{December 16, 2022}


A few afternoons ago, I was having tea with some of the guests here at the monastery during the teatime Q\&A, and one of the topics that came up in our discussion was how to use certain teachings that we don't completely understand yet. How do we hold them when we're unsure about and haven't really experienced them ourselves or know what to think about them in relation to the teachings, the training? How do we use teachings as contemplative strategies rather than doctrines to believe in?

I thought that would be an interesting thing to explore a little bit tonight because it's a fundamental aspect of our practice in the way that the Buddha set it out. We aren't instructed to pick up teachings in a way that's just on a scholarly or intellectual level in order to present our views and opinions as being right and everybody else's view as being wrong. The Buddha pointed out the unskillfulness of that and how it doesn't actually get us anywhere other than ensnared and entangled with views, opinions, and conflict.

So how do we use teachings for the purpose of ultimately realizing the end of suffering? Using them as contemplative strategies is generally how he instructs us to hold the practice, the teachings, and the discipline, both the monastic discipline and the lay precepts.

The word ``\emph{opanayiko}'' that we see in the chanting means ``to be brought in.'' The teachings are to be brought in, not just to be learned on an intellectual level. It doesn't mean that the Buddha doesn't have some categorical teachings that he presents as facts or as accurate representation of natural laws. For example, the teatime discussion that we were having was particularly around kamma and rebirth.

There are a number of teachings where he categorically says, ``As I understand it, as I've experienced it, this is how it works.'' The teachings in the category of Right View fall into that category: kamma and rebirth being two of them.The Four Noble Truths is also a categorical teaching. But even with categorical teachings, they're not to be picked up just for the sake of understanding on a conceptual level but are to be applied as a strategy for reaching the end of suffering.

%\looseness=-1
Both categorical teachings and other teachings that aren't categorical but are more reflective should be used as skillful means. As for the ones that we don't understand or haven't experienced for ourselves, we can pick them up and use them as working hypotheses. We haven't necessarily proven them to ourselves, but we can pick them up on faith, not blind faith, but with the thought, ``Okay, what I've picked up and applied before from the Buddha's teachings seem to have had an effect. They seem to work. The Buddha also taught \emph{this}, which I haven't really experienced. I'll pick it up on faith because a lot of everything else that he taught has worked so far, and I'll try it out. I'll pick it up as a hypothesis, try it out, see what happens. What are the effects in my mind, \mbox{in my heart?''}

%\looseness=-1
It's a learning process. It's an educational process that we're picking up. And we're using it with a spirit of inquiry, not in the spirit of doctrinal belief. One of the best examples of this that sometimes gets misunderstood, particularly when one first encounters the teachings, is what sometimes is referred to as the doctrine of not-self, \emph{anatta}, or the characteristic of not-self. Just the way it's phrased sometimes, for example using the word ``doctrine'' or ``characteristic,'' starts to turn it into something that is more along the lines of a description of a belief, rather than as a contemplative strategy. And I think it's really useful to undo that habit because a lot of people who are trying to understand this teaching end up with questions such as, ``Well, the Buddha says there is no self. If there is no self, what is all this: this body, this mind?'' And then from that comes a whole barrage of questions that just run in circles and don't lead to any kind \mbox{of deep understanding.}

\enlargethispage{1\baselineskip}
I think about a time quite a few years ago when I was staying at Pacific Hermitage in White Salmon, Washington, and Ajahn Sudanto and I were invited to give a presentation at a local community college on the other side of the river in Oregon. They were studying the world's religions and asked us to come and share some of our understanding of Buddhist practice. It was a great class, very lively people, with very intelligent, inquisitive minds. Eventually it came around to this particular teaching on not-self, and one young woman was asking a lot of questions about it. I explained how the teaching on not-self could be used as a contemplative strategy rather than as a doctrine of belief, as a way of deconstructing our sense of self.

I was trying to point out that the sense of self is merely a perception, but that we can run into problems when we say, ``There is no self.'' If we start categorically saying things like ``there is no self,'' then it can lead to confusion because the sense of self is such a strong perception, which most of us experience. It really feels like there's me here. I'm here. I'm this body. I'm this mind. Or, I am in control of this body or in control of this mind. The Buddha didn't declare, ``There is no self.'' There's a line in our morning chanting, ``\emph{sabbe dhammā anattā}'', and it's translated there as, ``There is no self in the created or the uncreated,'' which might lead to an annihilationist viewpoint. ``\emph{Sabbe dhammā anattā''} might be better translated as, ``All \emph{dhammas}, or all phenomena, are not self.'' We can see there's a difference between saying ``there is no self'' and ``all \emph{dhammas} are not self.'' The second translation refers more to that contemplative strategy.

During this discussion, I was explaining to her that as a contemplative strategy, we can go through what we normally think of as self. This is laid out in very succinct phrases in the \emph{Anattalakkhaṇa Sutta} as the five aggregates: form, feeling, perception, mental formations, and consciousness. In some form or combination of those aggregates, singularly or collectively, we impute a sense of self onto them. We take \linebreak this body and/or this mind in its various forms to be who we are. The \emph{Anattalakkhaṇa Sutta} goes through each one of the aggregates and starts to unravel that perception with the Buddha's skillful \mbox{questioning of the bhikkhus.}

I described how we go through each of these aggregates. We look at the body. We investigate it, and we see it very clearly through a pristine lens of \emph{samādhi}, a collected mind that's been settled in meditation. We start to examine our assumptions of this body being who we are or something that we possess. We carefully examine it piece by piece, part by part, function by function, and ask ourselves, ``Is this something that I have ultimate control over? Can I say that this is who I am or something that I possess?''

We continue that investigation through the aspects of mind, including our feelings, the pleasantness or unpleasantness of experience, and the perceptions that we create. Perceptions are patterns that try to make sense of the world around us, based on data that comes in through all of the different sense bases: eye, ear, nose, tongue, body, and mind. We create mental perceptions to help give meaning to our experience.

\enlargethispage{1\baselineskip}
\looseness=-1
We ask the questions: ``Are they who we really are? Are they completely under our control? Are they unchanging?'' And so on with thoughts, mental formations, and sensory consciousness.

I was explaining this contemplation of going through them one by one with a clear mind and examining these aggregates to see whether they stand up to the test of something that we could call ``me'' or ``mine.'' The process is one of deconstructing our patterns of perception that impute a sense of self onto these five aggregates. It's a process of trying to uncover that this \mbox{isn't who I am.}

She was really intrigued by this and was processing out loud. She said, ``So, I think I get this. You go through this process of deconstructing who it is that you \emph{think} you are but really are not, as a contemplative strategy. And when you finally penetrate and realize that you really aren't those things, then what appears is who you \emph{really} are.''

And I paused for a moment as I was thinking, and then my response was, ``Well, not exactly. You go through this process of deconstructing the sense of self in relation to forms, feelings, perceptions, mental formations, and consciousness. The result is finding out who you really \emph{aren't}. There is no second stage of finding out who you really are.''

\looseness=-1
It's just that one has completely discovered who one really is not. And then leaves it there to see what happens in the mind. This is using the teaching on not-self as a contemplative strategy.

As the teachings of the Buddha point out, with the complete understanding through contemplation, we develop a healthy sense of disenchantment with all of these things that we've formerly assumed to be who we are. We don't look for a sense of security, a sense of stability in these aggregates that are constantly changing. That's what then leads us to relinquishment, to letting go of constantly moving within those realms of trying to create happiness, stability, permanence, and security in the world of the aggregates. That offers the opportunity to open up to what the Buddha described as the Unconditioned, that which does not participate in or partake of the world of the constructed, the fabricated. However, he doesn't go into redefining it and reifying it or concretizing any kind of experience into a sense of self.

\looseness=-1
There are many strategies similar to this particular contemplation that the Buddha gives to us; many teachings, many suttas that describe this process. There's a particular one that I've always liked. It's a very straightforward, simple simile where the Buddha is with a group of his disciples in Jeta's Grove, Anāthapiṇḍika's Park, a place where he spent a lot of time teaching.

He says, and I paraphrase, ``In this Jeta's Grove, if people were going around and collecting all the grass, leaves, twigs, and branches scattered around, and starting to walk off with them to take them to use them for something or to burn them, would you say, `Stop, stop, don't take any more of those twigs! Don't take us! Don't take me! Those twigs and grass and leaves, don't take me away! Don't burn me!' \mbox{Would you say that?''}


And the monks respond, ``Oh, of course not. And the Buddha says, \mbox{``Why is that?''}

``Because that doesn't have anything to do with us. That's not who we are. Those are just products of the forest.''

\looseness=-1
He said, ``Well, in the same way, you should not assume to take possession of, take hold of, to identify with that \mbox{which isn't yours.}

And what is it that isn't yours? Form, this body, isn't yours. Abandon it. Feeling isn't yours. Abandon it. Perceptions are not yours. Abandon them. Mental formations aren't yours. Abandon them. Consciousness isn't yours. Abandon it. In the same way that you abandon and don't take possession of and don't concern yourself with the forest products being taken away, don't regard any of these five aggregates as yourself. Abandon them. Abandon that which isn't yours.''

We can use this as a contemplative strategy to think, ``In the same way that all these bits and pieces of the forest are just mere elements, earth element in particular, we can regard the body as just mere elements: the organs, the fluids, the winds, the different aspects of heat, the four physical elements of this body. That's all it is. It has nothing to \mbox{do with me or mine.''}

There are many other suttas that work along similar lines to help deconstruct the sense of self. Pick them up, not as a doctrine, but as something to investigate.

Another teaching that I found really useful is something that Ajahn Sucitto has talked about in reference to the body and the mind. He talks about them in terms of intelligence. Normally, when we think of intelligence, we think of it in regard to the thinking mind that is involved with conceptualization, abstraction, and planning---very active mental activities. Often when we think of an intelligent person, he or she is someone who's able to process information and make associations well. It's a very useful tool that we have, this intellect and the intelligence within it. It helps us to organize data and store data in our memories so that if it's useful data, we can then retrieve it and help to organize our world in useful ways.

If it's not used skillfully, it can also be co-opted very easily by the defilements of desire, aversion, and delusion. It can be put to harmful ends, creating all sorts of stories and unskillful states of mind, and then we dwell and proliferate on those. Whether that intellectual ability is used to our benefit or not depends on what's in charge at any one time. It can be a very powerful tool.

There are also other kinds of intelligence that we don't really consider very much, such as bodily intelligence. We can observe the body as a very independent organism that has its own innate intelligence, which we don't pay much attention to. We generally go on automatic pilot with our body. Say you're sitting down reading a book, and all of a sudden you get hungry. You put the book down, get up, and go to the fridge. You pull out a bagel. You cut it in half without cutting your finger off, pop it into the toaster, turn around, go find some mayonnaise or mustard, slather it on the bagel, add some cheese, some sprouts, put it all together, take it back to the table, and sit down.

\looseness=-1
You're probably not very conscious of the activities that you're doing. You may be talking to a friend. You've got your earbuds in. You're talking to somebody over the phone, listening to some music, or listening to a Dhamma talk on not-self while you're making your sandwich. You've gone through this whole process, and you haven't had to think about it consciously. You haven't had to instruct your body how to get up, lift its foot, and move it forward, and then do the same with the other foot to go in the direction you want; how to open the fridge; how to cut the bagel; how to construct a sandwich. You're probably doing this without much awareness at all. The body is just doing all these activities that require incredible dexterity and precision. Meanwhile, you're talking to your friend and not even paying any attention to what you're doing. However, there's a general, vague awareness of motion, of something happening.

\looseness=-1
This is the bodily intelligence. The body has its own innate quality that helps it function, whether we're aware of it or not, whether we're conscious or not. It even functions when we're completely unconscious, like when under anesthesia during surgery; the body keeps on functioning. The heart keeps beating. The lungs keep breathing. The metabolic processes keep happening. When we are conscious, the body's movements and activities are done without much awareness. This is the innate bodily intelligence, and it's quite profound.

This is similar with the heart/mind---not the intellect, but the \emph{citta}. \emph{Citta} is a process with its own innate intelligence that feels, that is sensitive. It intuitively picks up and experiences perceptions and feelings.

For example, something comes along that's a moving experience. We see somebody who is having difficulty, struggling or suffering in some way, and it doesn't take any kind of thinking process to feel a response. The citta can automatically feel compassion, a sense of caring and concern. It's an intelligence of the citta, of the heart. Rational prompting is unnecessary to understand that kind of experience.

Or we do some sort of unskillful action that causes difficulty or harm, something we say or do that results in pain. If it's seen clearly and recognized, the heart responds automatically with a sense of, ``Oh God, why did I do that? That's so painful to see. I feel really bad about that.'' We don't have to go through a big intellectual process of analyzing step-by-step what happened. If we have cultivated a sense of conscience, this sense of skillful regret arises because the heart is impacted. This is the innate intelligence that our citta operates with, for better or for worse.

There are these kinds of intelligence that if we really bring them into regular contemplation and experience them, we can start to realize, ``Oh wow, my thoughts, my feelings, my perceptions, even my bodily activities keep on operating, whether I am conscious of them or not, whether I am in control or not, whether I'm present or not.'' And we can start to step back a little bit from this quality of our mind enmeshed with the object, where we identify so strongly with the body and the activities of the mind, even the very subtle fluctuations of the mind. All of these are happening on their own steam, without the presence of an agent that is in control. When we can use our meditation to quiet the mind and then start to examine these experiences, not thinking about them so much, but watching them as they manifest, with a little bit of growing distance between the quality of pure knowing and the activity itself, then this can give us a sense of this perception of not-self. With this perception, there can then be a little bit of disentanglement, deconstruction from this marriage to a sense of ``I am.''

\looseness=-1
I have found that these are useful tools. I have a lot of appreciation for the people who have articulated the teachings in this way. It brings the practice out of the realm of doctrinal belief and into the realm of a strategy that we can use, and that results in deep insight such that our perspective of who we are in this world starts to change. We discover ways to help free the heart from the entanglement that we've woven for ourselves.

%\enlargethispage{1\baselineskip}
When we do that, as I mentioned before, there arises a sense of healthy disenchantment. We can clearly see the ridiculousness of how we keep creating all of these paradigms in a way that we think will give us some sort of long-term gratification but doesn't. It's short term. Instead, we adopt new strategies to help steer the mind towards that clear seeing, an awareness of how much we construct this perception of self to provide a sense of security for ourselves. We don't need to do that. As we are able to realize, understand, and experience that deeply, then a profound sense of freedom can arise, an appreciation of letting go and relinquishing these old patterns. This can serve as the stepping stone into what is referred to as the entry into emptiness, the experience of emptiness of self.

\looseness=-1
There is a freedom that results from that, and we are able to use that as our point of leaping off into an experience of the Unconstructed, the Non-Manifest. We can directly realize the Four Noble Truths: understanding suffering, understanding how it's caused, and practicing the way leading to the complete ending of suffering. That's the potential from using all of these teachings as contemplative strategies rather than just as interesting doctrines to believe in.
